% Generated draft: requires mathematical and source review.
\documentclass[11pt]{article}
\usepackage[margin=25mm]{geometry}
\usepackage{fontspec}
\setmainfont{Latin Modern Roman}
\usepackage{amsmath,amssymb,mathtools,unicode-math}
\setmathfont{Latin Modern Math}
\usepackage{enumitem,needspace,xurl,hyperref}
\setlength{\parindent}{0pt}
\setlength{\parskip}{5pt}
\setlength{\emergencystretch}{3em}
\newcommand{\N}{\mathbb{N}}
\newcommand{\Z}{\mathbb{Z}}
\newcommand{\Q}{\mathbb{Q}}
\newcommand{\R}{\mathbb{R}}
\newcommand{\C}{\mathbb{C}}
\newcommand{\F}{\mathbb{F}}
\newcommand{\A}{\mathbb{A}}
\newcommand{\PP}{\mathbb{P}}
\newcommand{\E}{\mathbb{E}}
\newcommand{\eps}{\varepsilon}
\newcommand{\dd}{\,\mathrm{d}}
\begin{document}
\section*{UnsolvedMath ID: 10900003}
\textbf{Source problem code:} catalogue entry
\subsection{Definitions and mathematical statement}
Let \(M\) be a closed hyperbolic three-manifold, meaning a compact manifold without boundary admitting a complete metric of constant sectional curvature \(-1\). In the standard real-projective interpretation, a projective structure \(\sigma\) on \(M\) is an atlas into \(\R P^3\) with transition maps locally restricted from \(\operatorname{PGL}_4(\R)\). It has a developing map and holonomy
\[\operatorname{dev}_\sigma:\widetilde M\to\R P^3,\qquad
\rho_\sigma:\pi_1(M)\to\operatorname{PGL}_4(\R),\]
with \(\operatorname{dev}_\sigma(\gamma x)=\rho_\sigma(\gamma)\operatorname{dev}_\sigma(x)\).

Cooper asks whether every such structure is convex. Under the developing-map convention, this means that \(\operatorname{dev}_\sigma\) identifies \(\widetilde M\) with a convex domain \(\Omega\subset\R P^3\), so that the structure is the quotient of \(\Omega\) by its holonomy action. In an affine-chart version, \(\Omega\) lies in an affine chart and is convex in its affine coordinates. Proper convexity additionally requires the appropriate bounded-domain/projective-closure condition; the excerpt does not specify whether ``convex'' is intended to include that stronger condition.

The quantifiers are
\[\forall M\text{ closed hyperbolic of dimension }3,\quad
\forall\sigma\text{ projective on }M,\quad\sigma\text{ is convex}.\]
The structures range over all projective structures, without a restriction to the hyperbolic structure's deformation component or a preselected holonomy representation. Existence of a hyperbolic metric on \(M\) does not itself define \(\sigma\). The source's real-versus-other projective convention and exact convexity convention require confirmation; no extra properness or holonomy hypothesis is silently added.
\subsection{Short English statement}
Must every projective structure on a closed hyperbolic three-manifold develop injectively onto a convex domain? The question concerns all structures, with the source's exact convexity convention still to be fixed.
\subsection{Sources}
\begin{itemize}
\item Local UnsolvedMath registry, permanent ID \(10900003\), exact source code \texttt{catalogue entry}.
\item Cooper, attribution in Delp, Hoffoss and Manning, \emph{Problems In Groups, Geometry, and Three-Manifolds} (2015), subsection 1.3, PDF page 1: \url{https://arxiv.org/pdf/1512.04620#page=1}; abstract record \url{https://arxiv.org/abs/1512.04620}. Registry access 2026-07-29; source-tex extraction recorded as checked against rendered PDF text.
\end{itemize}
\end{document}
